% GENERATED FILE. Edit canonical lessons, then run npm run generate:books.
% canonical-source-hash: fnv1a64:df1573728f49f289
% canonical-lessons: ZH-C540-guguai, ZH-C540-lianlei, ZH-C540-nanmian, ZH-C540-qianyinhouguo, ZH-C540-zenmehuishi

\chapter{So That's Why}
\label{ch:zh-so-that-s-why}
\hlchaptermodality{540}

\begin{chapteropening}
\textbf{By the end of this chapter:} \emph{I can say someone is odd, that something drags others in, that something is bound to happen, tell the whole story and ask what is going on.}

Complete the last lesson of chapter 540: I can say someone is odd, that something drags others in, that something is bound to happen, tell the whole story and ask what is going on.
\end{chapteropening}

\section[\texorpdfstring{\zh{古怪} (gǔguài)}{古怪 (gǔguài)}]{\zh{古怪} (gǔguài) — eccentric; peculiar (of a person)}

\label{lesson:ZH-C540-guguai}



\begin{quote}
Before the new one: say the Mandarin for to come off worse, then the Mandarin for to get into trouble.
\end{quote}

\subsection*{You'll want to know: \zh{古怪}}
\textbf{\zh{古怪}} — \emph{gǔguài} — \textquotedblleft{}eccentric; peculiar (of a person)\textquotedblright{}.

\begin{quote}
\textbf{\zh{原来他是一个古怪的人}} \emph{yuánlái tā shì yí ge gǔguài de rén} — So he's an odd sort of man.
\end{quote}

\begin{cousinweb}
\textbf{\zh{古怪}} (\emph{gǔguài}) is built from \zh{古} 'old' and \zh{怪} 'strange', both inherited from Middle Chinese, \emph{kuX} and \emph{kweajH} (Guangyun, in Baxter's transcription), which glosses the first 'of old' and the second 'strange'.
\end{cousinweb}

\begin{grammarlens}[title={What you've built}]
One more word for saying what would happen.
\end{grammarlens}

\subsection*{Your turn}
\begin{itemize}
\raggedright
  \item \emph{Say it:} \emph{gǔguài}
  \item \emph{Say it:} \emph{gǔguài}, once more
  \item \emph{Say it:} \emph{yuánlái tā shì yí ge gǔguài de rén} — So he's an odd sort of man.
  \item \emph{Recall:} say the Mandarin for to come off worse, then the Mandarin for to get into trouble, then say \emph{gǔguài} again
\end{itemize}

\begin{tcolorbox}[breakable,colback=teal!4,colframe=teal!35!black,title={Before you move on}]
What does \zh{古怪} mean? (\textquotedblleft{}Eccentric; peculiar (of a person)\textquotedblright{}.) Say it once more.
\end{tcolorbox}
\section[liánlei]{liánlei — to drag someone into trouble}

\label{lesson:ZH-C540-lianlei}



\begin{quote}
Before the new one: say the Mandarin for to get into trouble, then the Mandarin for eccentric.
\end{quote}

\subsection*{You'll want to know: liánlei}
\emph{liánlei} — \textquotedblleft{}to drag someone into trouble\textquotedblright{}.

\begin{quote}
\emph{yàoshi wǒ chūshì le jiù huì liánlei dàjiā} — If anything happens to me, everyone will be dragged into it.
\end{quote}

Its characters come later; for now it is a word for the ear.

\begin{cousinweb}
\emph{liánlei} is built from \emph{lián} 'to join' and \emph{lei}, both inherited from Middle Chinese, \emph{ljen} and \emph{ljweH} (Guangyun, in Baxter's transcription), which glosses the first 'to join, to continue' and the second as being caught up in another's offence. The second is said light here.
\end{cousinweb}

\begin{grammarlens}[title={What you've built}]
One more word for saying what would happen.
\end{grammarlens}

\subsection*{Your turn}
\begin{itemize}
\raggedright
  \item \emph{Say it:} \emph{liánlei}
  \item \emph{Say it:} \emph{liánlei}, once more
  \item \emph{Say it:} \emph{yàoshi wǒ chūshì le jiù huì liánlei dàjiā} — If anything happens to me, everyone will be dragged into it.
  \item \emph{Recall:} say the Mandarin for to get into trouble, then the Mandarin for eccentric, then say \emph{liánlei} again
\end{itemize}

\begin{tcolorbox}[breakable,colback=teal!4,colframe=teal!35!black,title={Before you move on}]
What does liánlei mean? (\textquotedblleft{}To drag someone into trouble\textquotedblright{}.) Say it once more.
\end{tcolorbox}
\section[nánmiǎn]{nánmiǎn — bound to happen; hard to avoid}

\label{lesson:ZH-C540-nanmian}



\begin{quote}
Before the new one: say the Mandarin for eccentric, then the Mandarin for to drag someone into trouble.
\end{quote}

\subsection*{You'll want to know: nánmiǎn}
\emph{nánmiǎn} — \textquotedblleft{}bound to happen; hard to avoid\textquotedblright{}.

\begin{quote}
\emph{rúguǒ shì dì yī cì zuò nánmiǎn huì chūcuò} — If it's your first time, mistakes are bound to happen.
\end{quote}

Its characters come later; for now it is a word for the ear.

\begin{cousinweb}
\emph{nánmiǎn} is built from \zh{难} 'hard' and \emph{miǎn} 'to escape', both inherited from Middle Chinese, \emph{nan} and \emph{mjenX} (Guangyun, in Baxter's transcription), which glosses the first 'difficult' and the second 'to stop' and 'to get free of': hard to get free of.
\end{cousinweb}

\begin{grammarlens}[title={What you've built}]
One more word for saying what would happen.
\end{grammarlens}

\subsection*{Your turn}
\begin{itemize}
\raggedright
  \item \emph{Say it:} \emph{nánmiǎn}
  \item \emph{Say it:} \emph{nánmiǎn}, once more
  \item \emph{Say it:} \emph{rúguǒ shì dì yī cì zuò nánmiǎn huì chūcuò} — If it's your first time, mistakes are bound to happen.
  \item \emph{Recall:} say the Mandarin for eccentric, then the Mandarin for to drag someone into trouble, then say \emph{nánmiǎn} again
\end{itemize}

\begin{tcolorbox}[breakable,colback=teal!4,colframe=teal!35!black,title={Before you move on}]
What does nánmiǎn mean? (\textquotedblleft{}Bound to happen; hard to avoid\textquotedblright{}.) Say it once more.
\end{tcolorbox}
\section[\texorpdfstring{\zh{前因后果} (qiányīn hòuguǒ)}{前因后果 (qiányīn hòuguǒ)}]{\zh{前因后果} (qiányīn hòuguǒ) — the whole chain of causes and effects}

\label{lesson:ZH-C540-qianyinhouguo}



\begin{quote}
Before the new one: say the Mandarin for to drag someone into trouble, then the Mandarin for bound to happen.
\end{quote}

\subsection*{You'll want to know: \zh{前因后果}}
\textbf{\zh{前因后果}} — \emph{qiányīn hòuguǒ} — \textquotedblleft{}the whole chain of causes and effects\textquotedblright{}.

\begin{quote}
\emph{yàoshi nǐ zhīdào qiányīn hòuguǒ jiù bú huì guài tā le} — If you knew the whole story, you wouldn't blame him.
\end{quote}

\begin{cousinweb}
\textbf{\zh{前因后果}} (\emph{qiányīn hòuguǒ}) is built from \zh{前} 'before', \zh{因} 'a cause', \zh{后} 'after' and \zh{果} 'fruit', all inherited from Middle Chinese, \emph{dzen}, \emph{'jin}, \emph{huwX} and \emph{kwaX} (Guangyun, in Baxter's transcription): the cause that came before and the fruit that came after.
\end{cousinweb}

\begin{grammarlens}[title={What you've built}]
One more word for saying what would happen.
\end{grammarlens}

\subsection*{Your turn}
\begin{itemize}
\raggedright
  \item \emph{Say it:} \emph{qiányīn hòuguǒ}
  \item \emph{Say it:} \emph{qiányīn hòuguǒ}, once more
  \item \emph{Say it:} \emph{yàoshi nǐ zhīdào qiányīn hòuguǒ jiù bú huì guài tā le} — If you knew the whole story, you wouldn't blame him.
  \item \emph{Recall:} say the Mandarin for to drag someone into trouble, then the Mandarin for bound to happen, then say \emph{qiányīn hòuguǒ} again
\end{itemize}

\begin{tcolorbox}[breakable,colback=teal!4,colframe=teal!35!black,title={Before you move on}]
What does \zh{前因后果} mean? (\textquotedblleft{}The whole chain of causes and effects\textquotedblright{}.) Say it once more.
\end{tcolorbox}
\section[\texorpdfstring{\zh{怎么回事} (zěnme huí shì)}{怎么回事 (zěnme huí shì)}]{\zh{怎么回事} (zěnme huí shì) — what is going on here; what's the matter}

\label{lesson:ZH-C540-zenmehuishi}



\begin{quote}
Before the new one: say the Mandarin for bound to happen, then the Mandarin for the whole chain of causes.
\end{quote}

\subsection*{You'll want to know: \zh{怎么回事}}
\textbf{\zh{怎么回事}} — \emph{zěnme huí shì} — \textquotedblleft{}what is going on here; what's the matter\textquotedblright{}.

\begin{quote}
\textbf{\zh{这是怎么回事}} \emph{zhè shì zěnme huí shì} — What's going on here?
\end{quote}

\begin{cousinweb}
\textbf{\zh{怎么回事}} (\emph{zěnme huí shì}) is built from \zh{怎么} 'how', \zh{回} 'a time' and \zh{事} 'a matter'. \zh{回} and \zh{事} are inherited from Middle Chinese, \emph{hwoj} and \emph{dzriH} (Guangyun, in Baxter's transcription); no dictionary we use records an older form of \zh{怎}, so its history is not settled.
\end{cousinweb}

\begin{grammarlens}[title={What you've built}]
One more word for saying what would happen.
\end{grammarlens}

\subsection*{Your turn}
\begin{itemize}
\raggedright
  \item \emph{Say it:} \emph{zěnme huí shì}
  \item \emph{Say it:} \emph{zěnme huí shì}, once more
  \item \emph{Say it:} \emph{zhè shì zěnme huí shì} — What's going on here?
  \item \emph{Recall:} say the Mandarin for bound to happen, then the Mandarin for the whole chain of causes, then say \emph{zěnme huí shì} again
\end{itemize}

\begin{tcolorbox}[breakable,colback=teal!4,colframe=teal!35!black,title={Before you move on}]
What does \zh{怎么回事} mean? (\textquotedblleft{}What is going on here; what's the matter\textquotedblright{}.) Say it once more.
\end{tcolorbox}
